% !TeX program = xelatex
% ============================================================
%  CHƯƠNG 3: THIẾT KẾ VÀ TRIỂN KHAI TESTBED (PHASE 1)
%  PAD-ONAP: Proactive AI-Driven DDoS Defense on ONAP
% ============================================================
\documentclass[12pt,a4paper]{article}

\usepackage[a4paper,margin=2.5cm]{geometry}
\usepackage{fontspec}
\usepackage{polyglossia}
\setmainlanguage{vietnamese}
\setotherlanguage{english}
\setmainfont{Times New Roman}

\usepackage{amsmath,amssymb}
\usepackage{graphicx}
\usepackage{booktabs}
\usepackage{longtable}
\usepackage{array}
\usepackage{float}
\usepackage{enumitem}
\usepackage{listings}
\usepackage{xcolor}
\usepackage{tikz}
\usepackage[colorlinks=true,linkcolor=blue,citecolor=blue,urlcolor=blue]{hyperref}

\setlength{\parindent}{1.2em}
\setlength{\parskip}{0.5em}

\lstset{
  basicstyle=\ttfamily\footnotesize,
  backgroundcolor=\color{gray!8},
  frame=single,
  breaklines=true,
  columns=fullflexible,
  keepspaces=true,
}

\title{\textbf{Chương 3: Thiết Kế và Triển Khai Testbed}\\[0.4em]
\large PAD-ONAP: Proactive AI-Driven DDoS Defense on ONAP\\[0.3em]
\normalsize Phase 1 --- Hạ Tầng Thu Thập Telemetry và Mô Phỏng Mạng}
\author{}
\date{}

\begin{document}
\maketitle
\tableofcontents
\newpage

% ============================================================
\section{Tổng Quan Kiến Trúc Testbed}
% ============================================================

Mục tiêu của Phase 1 là xây dựng một môi trường thử nghiệm (\textit{testbed}) có thể tái tạo chính xác luồng dữ liệu telemetry từ hạ tầng mạng thực tế, đồng thời cung cấp giao diện chuẩn cho khối AI ở Phase 2 tiêu thụ. Testbed được thiết kế theo nguyên tắc \textit{độc lập phần cứng} --- toàn bộ chạy dưới dạng container Docker trên một máy phổ thông, không yêu cầu thiết bị mạng vật lý.

\subsection{Vai trò trong toàn bộ pipeline PAD-ONAP}

PAD-ONAP tổ chức luồng xử lý thành bốn mô-đun (M1--M4) như Bảng~\ref{tab:pipeline_overview} mô tả. Phase 1 triển khai mô-đun M1.

\begin{table}[H]
\centering
\caption{Tổng quan bốn mô-đun trong PAD-ONAP}
\label{tab:pipeline_overview}
\begin{tabular}{clp{7cm}l}
\toprule
\textbf{Mô-đun} & \textbf{Tên} & \textbf{Chức năng} & \textbf{Phase} \\
\midrule
M1 & Streaming Telemetry  & Thu thập dữ liệu mạng theo thời gian thực (gNMI), đẩy lên Kafka & 1 \\
M2 & AI Detection \& Forecast & Phát hiện tấn công (XGBoost) và dự báo (Transformer+LSTM) & 2 \\
M3 & ONAP Orchestration & Ánh xạ điểm rủi ro AI sang chính sách VNF qua Policy Framework & 3 \\
M4 & NFV Enforcement    & Triển khai và mở rộng VNF scrubber/rate-limiter & 3 \\
\bottomrule
\end{tabular}
\end{table}

\subsection{Luồng dữ liệu Phase 1}

Luồng dữ liệu được thiết kế theo mô hình pipeline tuyến tính, mỗi thành phần có vai trò độc lập và giao tiếp qua các giao thức chuẩn:

\begin{center}
\small
\texttt{gNMI Simulator} $\xrightarrow{\text{HTTP/500ms}}$
\texttt{kafka\_producer} $\xrightarrow{\text{Kafka}}$
\texttt{pad.telemetry.raw} $\xrightarrow{\text{Kafka}}$
\texttt{flink\_processor} $\xrightarrow{\text{Kafka}}$
\texttt{pad.telemetry.features} $\xrightarrow{\text{Kafka}}$
\texttt{live\_pipeline (AI)}
\end{center}

\vspace{0.5em}
\noindent Ngoài luồng Kafka đầy đủ, hệ thống còn hỗ trợ chế độ polling HTTP đơn giản hơn thông qua \texttt{NetFlow Collector}, phục vụ mục đích kiểm thử nhanh mà không cần khởi động toàn bộ hạ tầng:

\begin{center}
\small
\texttt{gNMI Simulator} $\xrightarrow{\text{HTTP}}$
\texttt{NetFlow Collector} $\xrightarrow{\text{GET /flows/latest}}$
\texttt{live\_pipeline (AI)}
\end{center}

\subsection{Yêu cầu hệ thống}

Testbed được thiết kế chạy trên phần cứng phổ thông, không yêu cầu thiết bị chuyên dụng:

\begin{table}[H]
\centering
\caption{Yêu cầu phần cứng tối thiểu và cấu hình thực tế}
\label{tab:hardware}
\begin{tabular}{lll}
\toprule
\textbf{Tài nguyên} & \textbf{Yêu cầu tối thiểu} & \textbf{Cấu hình thực tế} \\
\midrule
CPU  & 4 cores & AMD Ryzen 7 (8 cores) \\
RAM  & 8 GB   & 20 GB \\
GPU  & Không bắt buộc & NVIDIA RTX 3050 Ti (4 GB VRAM) \\
Lưu trữ & 20 GB & SSD \\
OS   & Ubuntu 20.04+ / Windows 11 & Windows 11 + WSL2 \\
Docker & 24.0+ & Docker Desktop 25.0 \\
\bottomrule
\end{tabular}
\end{table}

% ============================================================
\section{Mô-đun gNMI Simulator}
% ============================================================

\subsection{Mục đích và thiết kế}

Trong hệ thống thực tế, dữ liệu telemetry được lấy từ các thiết bị mạng (router, switch) thông qua giao thức gNMI (\textit{gRPC Network Management Interface}) với mô hình YANG. Để tạo ra môi trường thử nghiệm độc lập phần cứng, một máy chủ gNMI mô phỏng (\textit{gNMI Simulator}) được triển khai dưới dạng REST API, thay thế thiết bị mạng vật lý.

Trình mô phỏng mô hình hóa ba thiết bị mạng \texttt{r1}, \texttt{r2}, \texttt{r3} với các chỉ số telemetry thay đổi liên tục. Mỗi thiết bị duy trì một vòng lặp nền cập nhật số liệu theo chu kỳ 1 giây dựa trên mô hình nhiễu Gaussian, đồng thời phản ánh kịch bản tấn công khi được kích hoạt.

\subsection{Chỉ số telemetry được mô phỏng}

Mỗi bản tin telemetry bao gồm các chỉ số mạng được liệt kê trong Bảng~\ref{tab:gnmi_metrics}. Các chỉ số này ánh xạ trực tiếp sang 17 đặc trưng đầu vào của mô hình AI (xem Mục~\ref{sec:features}).

\begin{table}[H]
\centering
\caption{Các chỉ số telemetry do gNMI Simulator tạo ra}
\label{tab:gnmi_metrics}
\begin{tabular}{llp{7cm}}
\toprule
\textbf{Chỉ số} & \textbf{Đơn vị} & \textbf{Ý nghĩa} \\
\midrule
\texttt{in\_pkts} / \texttt{out\_pkts}  & gói/s & Tốc độ gói đến/đi \\
\texttt{in\_bytes} / \texttt{out\_bytes} & byte/s & Băng thông đến/đi \\
\texttt{tcp\_ratio} & [0, 1] & Tỉ lệ lưu lượng TCP \\
\texttt{udp\_ratio} & [0, 1] & Tỉ lệ lưu lượng UDP \\
\texttt{icmp\_ratio} & [0, 1] & Tỉ lệ lưu lượng ICMP \\
\texttt{syn\_ratio} & [0, 1] & Tỉ lệ gói SYN trong TCP \\
\texttt{fin\_ratio} & [0, 1] & Tỉ lệ gói FIN trong TCP \\
\texttt{src\_ip\_entropy} & [0, 1.5] & Entropy kích thước gói nguồn (proxy) \\
\texttt{dst\_ip\_entropy} & [0, 1.5] & Entropy kích thước gói đích (proxy) \\
\texttt{src\_port\_entropy} & [0, 7] & Entropy cổng nguồn (Shannon) \\
\texttt{dst\_port\_entropy} & [0, 7] & Entropy cổng đích (Shannon) \\
\texttt{avg\_pkt\_size} & byte & Kích thước gói trung bình \\
\texttt{new\_flows\_rate} & flow/s & Tốc độ luồng mới xuất hiện \\
\texttt{flow\_duration\_mean} & s & Thời gian sống trung bình của luồng \\
\texttt{iat\_mean\_ms} & ms & Thời gian giữa gói tin (trung bình) \\
\texttt{iat\_std\_ms} & ms & Thời gian giữa gói tin (độ lệch chuẩn) \\
\bottomrule
\end{tabular}
\end{table}

\subsection{API endpoints}

\begin{table}[H]
\centering
\caption{REST API của gNMI Simulator}
\label{tab:gnmi_api}
\begin{tabular}{llp{7cm}}
\toprule
\textbf{Phương thức} & \textbf{Endpoint} & \textbf{Chức năng} \\
\midrule
GET  & \texttt{/metrics}               & Tổng hợp số liệu của tất cả thiết bị \\
GET  & \texttt{/metrics/\{device\}}    & Số liệu của thiết bị đơn lẻ (r1/r2/r3) \\
GET  & \texttt{/health}                & Kiểm tra trạng thái hoạt động \\
POST & \texttt{/attack/start}          & Kích hoạt kịch bản tấn công \\
POST & \texttt{/attack/stop}           & Dừng tấn công, trở về traffic bình thường \\
POST & \texttt{/attack/ramp}           & Tấn công leo thang từ từ (ramp-up) \\
\bottomrule
\end{tabular}
\end{table}

\subsection{Hiệu chỉnh đặc trưng (Feature Calibration)}

Một thách thức quan trọng trong quá trình triển khai là sự khác biệt về nghĩa ngữ nghĩa (\textit{semantic mismatch}) giữa dữ liệu gNMI và không gian đặc trưng mà scaler của mô hình AI đã được huấn luyện.

Cụ thể, trong pipeline training (Phase 2), \texttt{src\_ip\_entropy} không phải entropy Shannon của địa chỉ IP mà là entropy của histogram kích thước gói (\textit{Avg Fwd Segment Size $\times 10$ bins}), với phân phối huấn luyện: $\mu = 0.12$, $\sigma = 0.50$. Nếu gNMI Simulator đẩy trực tiếp giá trị entropy IP (thường trong khoảng 4--7), đặc trưng đó sẽ nằm 13$\sigma$ ngoài phân phối huấn luyện, khiến mô hình dự báo sai hoàn toàn.

Để khắc phục, bộ sinh dữ liệu tổng hợp (\texttt{SyntheticFlowGenerator}) được hiệu chỉnh theo bảng ánh xạ sau (Bảng~\ref{tab:calibration}):

\begin{table}[H]
\centering
\caption{Bảng hiệu chỉnh ánh xạ gNMI → đặc trưng huấn luyện}
\label{tab:calibration}
\begin{tabular}{p{4.5cm}p{5cm}p{4cm}}
\toprule
\textbf{Đặc trưng} & \textbf{Công thức hiệu chỉnh} & \textbf{Lý do} \\
\midrule
\texttt{src\_ip\_entropy}   & $0.55 \times (1 - \alpha)$ & Proxy pkt-size entropy \\
\texttt{dst\_ip\_entropy}   & $0.70 \times (1 - \alpha)$ & Proxy pkt-size entropy \\
\texttt{syn\_ratio}         & $\min(0.02,\ \text{gNMI} \times 0.10)$ & Scale ×0.1 \\
\texttt{new\_flows\_rate}   & $\min(200000,\ \text{gNMI} \times 600)$ & Scale ×600 \\
\texttt{inter\_arrival\_mean} & $\text{gNMI (ms)} \times 60$ & Scale ×60 \\
\texttt{flow\_duration\_mean} & $\text{gNMI} \times 11$ (bình thường), $\times 2$ (tấn công) & Theo cường độ \\
\bottomrule
\end{tabular}
\end{table}

\noindent trong đó $\alpha \in [0,1]$ là cường độ tấn công được tính từ:
\[
\alpha = \min\!\left(1,\; 0.5 \cdot \frac{\text{udp\_ratio} - 0.6}{0.38} + 0.4 \cdot \frac{\text{in\_pkts} - 20000}{980000} + 0.1 \cdot \frac{\text{syn\_ratio} - 0.4}{0.59}\right)
\]

% ============================================================
\section{Mô-đun NetFlow Collector}
% ============================================================

\subsection{Chức năng và chế độ hoạt động}

NetFlow Collector đảm nhận vai trò trung gian giữa dữ liệu telemetry thô (từ gNMI) và pipeline AI. Thành phần này hỗ trợ hai chế độ vận hành:

\begin{enumerate}[leftmargin=2em]
\item \textbf{Chế độ NetFlow thực (\texttt{--mode netflow})}: Lắng nghe gói tin NetFlow v5/v9 trên cổng UDP 6343, phân tích cú pháp nhị phân và tổng hợp thành vector đặc trưng 17 chiều. Chế độ này được dùng khi kết nối với thiết bị mạng vật lý có hỗ trợ NetFlow.

\item \textbf{Chế độ tổng hợp (\texttt{--mode synthetic})}: Định kỳ lấy dữ liệu từ gNMI Simulator qua HTTP (mặc định mỗi 1 giây) và áp dụng \texttt{SyntheticFlowGenerator} để tính vector đặc trưng. Chế độ này được dùng trong testbed, cho phép kiểm thử toàn bộ pipeline mà không cần thiết bị vật lý.
\end{enumerate}

\subsection{FlowFeatureExtractor và bộ đặc trưng 17 chiều}
\label{sec:features}

Từ cả hai chế độ, \texttt{FlowFeatureExtractor} tạo ra vector đặc trưng chuẩn hóa gồm 17 thành phần (Bảng~\ref{tab:features_groups}). Vector này có cấu trúc và thứ tự nhất quán tuyệt đối với đầu vào của \texttt{StandardScaler} và mô hình AI.

\begin{table}[H]
\centering
\caption{Bộ đặc trưng 17 chiều và nhóm đặc trưng}
\label{tab:features_groups}
\begin{tabular}{llp{7.5cm}}
\toprule
\textbf{Nhóm} & \textbf{Đặc trưng} & \textbf{Mô tả} \\
\midrule
\multirow{4}{*}{Tốc độ}
  & \texttt{pkt\_rate}           & Tổng gói/giây (vào + ra) \\
  & \texttt{byte\_rate}          & Tổng byte/giây (vào + ra) \\
  & \texttt{new\_flows\_rate}    & Số luồng mới xuất hiện/giây \\
  & \texttt{flow\_duration\_mean}& Thời gian trung bình một luồng (giây) \\
\midrule
\multirow{4}{*}{Entropy}
  & \texttt{src\_ip\_entropy}    & Entropy kích thước gói nguồn (proxy) \\
  & \texttt{dst\_ip\_entropy}    & Entropy kích thước gói đích (proxy) \\
  & \texttt{src\_port\_entropy}  & Entropy Shannon cổng nguồn \\
  & \texttt{dst\_port\_entropy}  & Entropy Shannon cổng đích \\
\midrule
\multirow{5}{*}{Giao thức}
  & \texttt{proto\_dist\_tcp}    & Tỉ lệ lưu lượng TCP \\
  & \texttt{proto\_dist\_udp}    & Tỉ lệ lưu lượng UDP \\
  & \texttt{proto\_dist\_icmp}   & Tỉ lệ lưu lượng ICMP \\
  & \texttt{syn\_ratio}          & Tỉ lệ gói SYN \\
  & \texttt{fin\_ratio}          & Tỉ lệ gói FIN \\
\midrule
\multirow{2}{*}{Kích thước gói}
  & \texttt{avg\_pkt\_size}      & Kích thước gói trung bình (byte) \\
  & \texttt{pkt\_size\_std}      & Độ lệch chuẩn kích thước gói \\
\midrule
\multirow{2}{*}{Inter-arrival}
  & \texttt{inter\_arrival\_mean}& Khoảng thời gian giữa gói (trung bình, ms) \\
  & \texttt{inter\_arrival\_std} & Khoảng thời gian giữa gói (độ lệch chuẩn, ms) \\
\bottomrule
\end{tabular}
\end{table}

\subsection{REST API của Collector}

NetFlow Collector cung cấp HTTP API cho live pipeline tiêu thụ:

\begin{table}[H]
\centering
\caption{REST API của NetFlow Collector (cổng 7070)}
\label{tab:collector_api}
\begin{tabular}{llp{7.5cm}}
\toprule
\textbf{Phương thức} & \textbf{Endpoint} & \textbf{Phản hồi} \\
\midrule
GET & \texttt{/flows/latest} & Vector đặc trưng 17 chiều mới nhất dạng JSON \\
GET & \texttt{/flows}        & $N$ vector gần nhất (tham số \texttt{?n=}) \\
GET & \texttt{/health}       & Trạng thái hoạt động của collector \\
\midrule
\multicolumn{3}{l}{\textit{Phản hồi 204 No Content nếu chưa có dữ liệu}} \\
\bottomrule
\end{tabular}
\end{table}

% ============================================================
\section{Hạ Tầng Kafka và Sliding-Window Feature Extraction}
% ============================================================

\subsection{Lý do lựa chọn Kafka KRaft}

Trong hạ tầng thực tế của ONAP, Apache Kafka đóng vai trò message broker trung tâm cho luồng telemetry (DCAE/DMaaP). Để testbed phản ánh đúng kiến trúc thực tế, hệ thống triển khai một cụm Kafka đơn node theo chế độ KRaft (không cần Zookeeper), giảm overhead xuống còn khoảng 800~MB RAM so với 1.5~GB khi dùng Zookeeper.

\subsection{Cấu trúc topic Kafka}

\begin{table}[H]
\centering
\caption{Các Kafka topic trong PAD-ONAP}
\label{tab:kafka_topics}
\begin{tabular}{lll}
\toprule
\textbf{Topic} & \textbf{Nguồn} & \textbf{Người tiêu thụ} \\
\midrule
\texttt{pad.telemetry.raw}      & \texttt{kafka\_producer.py}  & \texttt{flink\_processor.py} \\
\texttt{pad.telemetry.features} & \texttt{flink\_processor.py} & \texttt{live\_pipeline.py} \\
\bottomrule
\end{tabular}
\end{table}

\subsection{kafka\_producer.py}

\texttt{kafka\_producer.py} poll gNMI Simulator theo chu kỳ (mặc định 0.5 giây), đóng gói số liệu thô thành bản tin JSON kèm timestamp ISO-8601 và đẩy lên \texttt{pad.telemetry.raw}.

Để đảm bảo độ ổn định khi triển khai lên server ONAP, component này được thiết kế với các cơ chế:

\begin{itemize}[leftmargin=2em]
\item \textbf{Exponential backoff khi kết nối}: Thời gian chờ bắt đầu từ 2 giây, tăng gấp đôi mỗi lần thất bại đến tối đa 60 giây.
\item \textbf{Gửi bất đồng bộ}: Sử dụng \texttt{producer.send()} với error callback thay vì blocking \texttt{future.get()}, tránh làm gián đoạn vòng lặp sản xuất.
\item \textbf{Tự động kết nối lại}: Khi phát hiện lỗi gửi, producer được tạo lại hoàn toàn thay vì giữ nguyên kết nối bị hỏng.
\item \textbf{Flush định kỳ}: Mỗi 5 giây thực hiện \texttt{producer.flush()} đảm bảo bản tin đến được broker.
\end{itemize}

\subsection{flink\_processor.py --- Sliding-Window Feature Aggregation}

\texttt{flink\_processor.py} triển khai bộ tổng hợp cửa sổ trượt (\textit{sliding-window aggregation}) trong Python thuần, không dùng JVM PyFlink (tiết kiệm $\sim$1.5~GB RAM). Tham số mặc định: cửa sổ 5 giây / slide 1 giây.

\begin{enumerate}[leftmargin=2em]
\item Consume bản tin từ \texttt{pad.telemetry.raw} bằng \texttt{KafkaConsumer} (non-blocking, \texttt{consumer\_timeout\_ms=300}).
\item Nạp mỗi snapshot vào hàng đợi \texttt{deque} có đánh timestamp.
\item Mỗi khi đến kỳ slide (1 giây), loại bỏ snapshot cũ hơn 5 giây, tính vector đặc trưng trung bình/phương sai trên cửa sổ hiện tại.
\item Publish kết quả lên \texttt{pad.telemetry.features} kèm metadata (window\_s, slide\_s, timestamp).
\end{enumerate}

\noindent Tương tự \texttt{kafka\_producer.py}, component này có cơ chế tự kết nối lại với exponential backoff và flush bắt buộc sau mỗi lần emit.

\subsection{Cấu hình Kafka dual-listener}

Để đồng thời phục vụ container Docker và Python chạy trực tiếp trên host (cần CUDA), Kafka được cấu hình hai listener:

\begin{table}[H]
\centering
\caption{Cấu hình dual-listener Kafka}
\label{tab:kafka_listener}
\begin{tabular}{llll}
\toprule
\textbf{Listener} & \textbf{Cổng} & \textbf{Địa chỉ quảng bá} & \textbf{Dùng cho} \\
\midrule
INTERNAL & 29092 & \texttt{kafka:29092} & Docker containers \\
EXTERNAL & 9092  & \texttt{<host-ip>:9092} & Python native (CUDA) \\
CONTROLLER & 9093 & \texttt{kafka:9093} & KRaft controller \\
\bottomrule
\end{tabular}
\end{table}

% ============================================================
\section{Hạ Tầng Docker Compose}
% ============================================================

\subsection{Cấu trúc services}

Toàn bộ hạ tầng được định nghĩa trong một file \texttt{docker-compose.yml} và có thể khởi động bằng một lệnh duy nhất. Bảng~\ref{tab:services} liệt kê tất cả services và tài nguyên phân bổ.

\begin{table}[H]
\centering
\caption{Docker Compose services và phân bổ tài nguyên}
\label{tab:services}
\begin{tabular}{lllr}
\toprule
\textbf{Service} & \textbf{Image} & \textbf{Cổng host} & \textbf{RAM tối đa} \\
\midrule
\texttt{kafka}            & apache/kafka:3.7.0        & 9092 (EXTERNAL) & 1~GB \\
\texttt{gnmi-simulator}   & Local Dockerfile          & 8080 & 256~MB \\
\texttt{netflow-collector}& python:3.11-slim          & 7070, 6343/udp & 256~MB \\
\texttt{prometheus}       & prom/prometheus:v2.51.0   & 9190 & 512~MB \\
\texttt{grafana}          & grafana/grafana:10.4.0    & 3001 & 256~MB \\
\texttt{metrics-exporter} & python:3.11-slim          & 9191 & 128~MB \\
\midrule
\multicolumn{3}{r}{\textbf{Tổng RAM}} & \textbf{$\approx$2.4~GB} \\
\bottomrule
\end{tabular}
\end{table}

\textbf{Lưu ý cổng}: Prometheus sử dụng cổng 9190 và Grafana sử dụng cổng 3001 thay vì mặc định 9090/3000, tránh xung đột với các service ONAP đang chạy trên server.

\subsection{Health check và restart policy}

Tất cả services đều có cấu hình health check và \texttt{restart: always} đảm bảo tự phục hồi khi gặp lỗi:

\begin{lstlisting}[language=yaml,caption={Ví dụ cấu hình health check của Kafka}]
healthcheck:
  test: ["/opt/kafka/bin/kafka-topics.sh",
         "--bootstrap-server", "localhost:9092", "--list"]
  interval: 10s
  timeout: 8s
  retries: 6
  start_period: 20s
restart: always
logging:
  driver: "json-file"
  options: {max-size: "20m", max-file: "5"}
\end{lstlisting}

\subsection{Cấu hình môi trường với file .env}

Mọi thông số phụ thuộc vào môi trường deployment được tập trung trong file \texttt{testbed/.env}:

\begin{table}[H]
\centering
\caption{Các biến môi trường chính trong \texttt{.env}}
\label{tab:env}
\begin{tabular}{llp{6cm}}
\toprule
\textbf{Biến} & \textbf{Mặc định} & \textbf{Mô tả} \\
\midrule
\texttt{PAD\_HOST}              & \texttt{localhost} & IP/hostname server ONAP \\
\texttt{PAD\_KAFKA\_PORT}       & 9092 & Cổng Kafka external \\
\texttt{PAD\_PROMETHEUS\_PORT}  & 9190 & Cổng Prometheus (tránh 9090 của ONAP) \\
\texttt{PAD\_GRAFANA\_PORT}     & 3001 & Cổng Grafana (tránh 3000 của ONAP) \\
\texttt{PAD\_MODEL\_DIR}        & \texttt{./pad\_onap\_v3/models} & Đường dẫn model AI \\
\bottomrule
\end{tabular}
\end{table}

% ============================================================
\section{Kịch Bản Tấn Công và Anomaly Injection}
% ============================================================

\subsection{Các kịch bản tấn công được hỗ trợ}

Module \texttt{anomaly\_injector} cung cấp 4 kịch bản tấn công DDoS có thể được kích hoạt thông qua REST API của gNMI Simulator:

\begin{table}[H]
\centering
\caption{Kịch bản tấn công trong testbed}
\label{tab:attacks}
\begin{tabular}{llp{8cm}}
\toprule
\textbf{Kịch bản} & \textbf{Loại} & \textbf{Tham số đặc trưng bị ảnh hưởng} \\
\midrule
UDP Flood       & Volumetric  & \texttt{udp\_ratio}$\uparrow$, \texttt{pkt\_rate}$\uparrow$, \texttt{dst\_port\_entropy}$\downarrow$ \\
SYN Flood       & Protocol    & \texttt{syn\_ratio}$\uparrow$, \texttt{tcp\_ratio}$\uparrow$, \texttt{fin\_ratio}$\downarrow$ \\
Bandwidth Ramp  & Leo thang   & Tăng dần \texttt{in\_pkts} theo thời gian \\
CPU Spike       & Resource    & \texttt{cpu\_pct}$\uparrow$, \texttt{queue\_depth\_pct}$\uparrow$ \\
\bottomrule
\end{tabular}
\end{table}

\subsection{Chuyển đổi sang phân phối huấn luyện}

Khi kích hoạt tấn công, \texttt{SyntheticFlowGenerator} điều chỉnh tất cả 17 đặc trưng theo cường độ tấn công $\alpha$, đảm bảo vector đặc trưng nằm trong phân phối mà mô hình AI đã học:

\begin{itemize}[leftmargin=2em]
\item Lưu lượng bình thường: $z$-score của mỗi đặc trưng nằm trong $[-1.5\sigma, +1.5\sigma]$.
\item Tấn công cường độ cao ($\alpha \to 1$): đặc trưng lệch rõ khỏi phân phối bình thường, kích hoạt phát hiện của XGBoost (confidence $> 0.95$) và trigger proactive của Transformer ($P_{30s} > 0.70$).
\end{itemize}

% ============================================================
\section{Live Pipeline --- Kết Nối Phase 1 và Phase 2}
% ============================================================

\subsection{Hai chế độ nguồn dữ liệu}

Module \texttt{live\_pipeline.py} đóng vai trò cầu nối cuối cùng, nhận vector đặc trưng từ Phase 1 và đưa vào InferenceEngine của Phase 2. Module hỗ trợ hai chế độ:

\begin{table}[H]
\centering
\caption{So sánh hai chế độ nguồn dữ liệu của live\_pipeline.py}
\label{tab:source_modes}
\begin{tabular}{lp{5cm}p{5cm}}
\toprule
 & \textbf{\texttt{--source http}} & \textbf{\texttt{--source kafka}} \\
\midrule
Độ phức tạp & Thấp (không cần Kafka) & Cao (cần full stack) \\
Mục đích    & Kiểm thử nhanh, debug & Production/demo với ONAP \\
Nguồn dữ liệu & GET /flows/latest từ Collector & Topic \texttt{pad.telemetry.features} \\
Độ trễ      & $\sim$1s (polling) & $\leq$1s (sliding window) \\
\bottomrule
\end{tabular}
\end{table}

\subsection{KafkaFeatureConsumer với auto-reconnect}

Trong chế độ Kafka, \texttt{KafkaFeatureConsumer} được thiết kế với khả năng tự phục hồi hoàn toàn:

\begin{itemize}[leftmargin=2em]
\item \textbf{Kết nối ban đầu}: Vòng lặp với exponential backoff, không giới hạn số lần thử.
\item \textbf{Poll không blocking}: \texttt{consumer\_timeout\_ms=500}, thu thập tất cả bản tin mới và chỉ giữ bản tin cuối cùng (tránh xử lý dữ liệu cũ tích lũy).
\item \textbf{Tự kết nối lại}: Khi phát hiện lỗi trong vòng lặp poll, đóng consumer hiện tại và tạo lại với logic backoff.
\end{itemize}

% ============================================================
\section{Scripts Triển Khai và Vận Hành}
% ============================================================

\subsection{Quy trình triển khai lên server ONAP}

\begin{lstlisting}[language=bash,caption={Triển khai PAD-ONAP Phase 1 lên server}]
# 1. Clone project
git clone <repo> Src_2 && cd Src_2

# 2. Setup một lần (kiểm tra port, tạo venv, cài deps)
chmod +x deploy/*.sh
./deploy/setup_server.sh

# 3. Chỉnh port nếu ONAP đang chiếm 9092/9090/3000
nano testbed/.env

# 4. Khởi động toàn bộ pipeline
./deploy/start.sh

# 5. Kiểm tra logs
tail -f logs/live_pipeline.log
\end{lstlisting}

\subsection{deploy/start.sh --- Orchestration script}

Script \texttt{deploy/start.sh} thực hiện tuần tự các bước:

\begin{enumerate}[leftmargin=2em]
\item Đọc cấu hình từ \texttt{testbed/.env}.
\item Tạo Python virtual environment và cài đặt dependencies nếu chưa có.
\item Khởi động Docker infrastructure (\texttt{kafka}, \texttt{gnmi-simulator}, \texttt{netflow-collector}).
\item Chờ Kafka đạt trạng thái \textit{healthy} (tối đa 60 giây).
\item Khởi động ba Python component native (nohup + PID tracking):
      \texttt{kafka\_producer.py}, \texttt{flink\_processor.py}, \texttt{live\_pipeline.py}.
\item In bảng trạng thái toàn bộ hệ thống và đường dẫn log.
\end{enumerate}

\subsection{Testbed Verification}

Script \texttt{scripts/verify\_testbed.sh} (327 dòng) kiểm tra 8 hạng mục tự động:

\begin{enumerate}[leftmargin=2em, itemsep=0pt]
\item System prerequisites (Python 3.10+, Docker $\geq$24.0, dependencies)
\item Project files (9 file bắt buộc)
\item gNMI Simulator (health, metrics, attack injection)
\item NetFlow Collector (health, /flows/latest, đủ 17 feature)
\item Anomaly Injector (4 kịch bản được định nghĩa)
\item Docker infrastructure (compose valid, image buildable)
\item Mininet Topology (syntax check)
\item Pipeline Integration (feature name matching giữa các module)
\end{enumerate}

% ============================================================
\section{Monitoring và Quan Sát}
% ============================================================

\subsection{Prometheus và Grafana}

Hệ thống monitoring bao gồm:

\begin{itemize}[leftmargin=2em]
\item \textbf{metrics-exporter}: Bridge từ gNMI Simulator sang Prometheus exposition format, cung cấp tất cả 17 chỉ số ở dạng time-series.
\item \textbf{Prometheus} (cổng 9190): Thu thập dữ liệu từ metrics-exporter với chu kỳ 5 giây, lưu trữ 7 ngày.
\item \textbf{Grafana} (cổng 3001): Dashboard trực quan hóa traffic theo thời gian thực, bao gồm tỉ lệ UDP/TCP, pkt\_rate, entropy, và trạng thái tấn công.
\end{itemize}

\subsection{JSONL inference log}

\texttt{live\_pipeline.py} ghi mỗi inference window ra file \texttt{logs/inference\_output.jsonl} theo định dạng JSONL (một JSON object mỗi dòng). Mỗi bản ghi bao gồm đầy đủ: kết quả detection, forecast 4 horizon, tier threat level, proactive trigger status và vector đặc trưng thô.

% ============================================================
\section{Kết Quả Kiểm Thử Phase 1}
% ============================================================

\subsection{Xác minh luồng end-to-end}

Kiểm thử end-to-end xác nhận luồng hoạt động chính xác từ gNMI Simulator đến InferenceEngine:

\begin{table}[H]
\centering
\caption{Kết quả kiểm thử Phase 1}
\label{tab:phase1_results}
\begin{tabular}{p{6cm}lp{5cm}}
\toprule
\textbf{Hạng mục kiểm thử} & \textbf{Kết quả} & \textbf{Ghi chú} \\
\midrule
gNMI Simulator khởi động & Đạt & Health check pass \\
NetFlow Collector synthetic mode & Đạt & 17 feature trả về đúng \\
Kafka producer gửi bản tin & Đạt & 50 msg/s ở interval 0.5s \\
Flink processor emit feature & Đạt & $\sim$1 vector/giây \\
Live pipeline (HTTP mode) & Đạt & $<$1ms round-trip \\
Live pipeline (Kafka mode) & Đạt & $<$1s end-to-end \\
Feature calibration (bình thường) & Đạt & $z$-score trong $\pm1.5\sigma$ \\
Attack injection → AI detect & Đạt & UDP\_Flood conf $>$0.95 \\
Proactive trigger P(t+30s) $>$0.70 & Đạt & Sau $\sim$12 cửa sổ \\
Docker restart recovery & Đạt & Kafka consumer tự kết nối lại \\
\bottomrule
\end{tabular}
\end{table}

\subsection{Latency đo được}

\begin{table}[H]
\centering
\caption{Latency thực đo của từng thành phần Phase 1}
\label{tab:latency}
\begin{tabular}{lrrr}
\toprule
\textbf{Thành phần} & \textbf{P50 (ms)} & \textbf{P99 (ms)} & \textbf{Mục tiêu} \\
\midrule
gNMI Simulator HTTP response  &  $<$1  &  $<$5  & $<$100 ms \\
kafka\_producer poll + send    & 0.5   & 2.1    & $<$100 ms \\
flink\_processor emit          & 1.0   & 4.5    & $<$100 ms \\
live\_pipeline inference (XGB) & 1.2   & 2.2    & $<$10 ms \\
live\_pipeline inference (TF)  & 5.1   & 7.7    & $<$10 ms \\
Tổng end-to-end (Kafka mode)   & $\sim$2000 & $\sim$2500 & $<$5000 ms \\
\bottomrule
\end{tabular}
\end{table}

% ============================================================
\section{Tóm Tắt Phase 1}
% ============================================================

Phase 1 đã xây dựng thành công hạ tầng testbed hoàn chỉnh, đáp ứng tất cả yêu cầu đặt ra:

\begin{itemize}[leftmargin=2em]
\item \textbf{Tái tạo trung thực} luồng telemetry từ mạng thực (gNMI $\to$ Kafka $\to$ Feature Vector).
\item \textbf{Hiệu chỉnh đặc trưng} giải quyết semantic mismatch giữa dữ liệu mô phỏng và không gian huấn luyện.
\item \textbf{Độ ổn định cao} với exponential backoff, auto-reconnect, health check và log rotation.
\item \textbf{Không xung đột} với ONAP thông qua cổng tách biệt (Prometheus 9190, Grafana 3001).
\item \textbf{Triển khai một lệnh} qua \texttt{./deploy/start.sh} trên bất kỳ server Linux nào có Docker.
\end{itemize}

\end{document}
